\documentclass[letterpaper,twocolumn,10pt]{article}
\usepackage{hyperref}
\usepackage{usenix}
\usepackage{booktabs}
\microtypesetup{spacing=false,kerning=false}

\begin{document}
\date{}
\title{\Large \bf Measuring Semantic Preservation in Passkey Migration}
\author{Anonymous Authors}
\maketitle

\begin{abstract}
An imported passkey can still authenticate while a property needed by an
application has changed. We present PasskeyTransit, a reproducible harness
that measures identity, WebAuthn behavior, extension state, and operational
guarantees separately. A balanced synthetic campaign exercised 6,144
credential--route--repetition attempts; 5,120 reached import and 1,024 were rejected
before a ceremony. Every executed assertion authenticated, yet 512 imported
attempts failed an executed browser behavior oracle for user verification or
largeBlob. These fractions describe the constructed matrix, not product
prevalence. Follow-up controls distinguish a project-authored CXP binding
from selected provider behavior. A Chromium virtual-authenticator transfer
preserved the signing credential but lost its PRF result; four public
assertions from a replay were independently reverified. Fenko authenticated
after CXF import while losing the original HMAC seeds. A Proton software
authenticator preserved both seeds and reproduced the expected PRF with user
verification. We provide claim boundaries and public verifiers. We did not
establish public CXP wire interoperability between providers or a verified
Android transfer.
\end{abstract}

\section{Introduction}
An accepted import and successful login are weak tests of passkey migration.
A relying party can also depend on the credential identifier, WebAuthn
extension output, user-verification behavior, and whether transfer retries
leave duplicate state. The Credential Exchange Format (CXF) represents
credentials~\cite{cxf}; the draft Credential Exchange Protocol (CXP)
addresses protected exchange~\cite{cxp}. WebAuthn authenticates an
assertion~\cite{webauthn}, but a valid signature does not by itself show
that every credential property survived the move.

We model a credential as $C=(I,K,F,S)$: identity, signing material,
observable functions, and operational guarantees. Our falsifiable
proposition is that valid authentication does not imply preservation of
every component. We contribute: (1) a claim matrix and property-specific
oracles; (2) a synthetic campaign with explicit denominators and
non-evaluable properties; (3) executable controls using two project-authored
runtimes, Chromium virtual authenticators, and selected open-source provider
code; and (4) an artifact whose public reports allow independent checks of
retained signatures and PRF digest relationships. This is a measurement
method with selected software cases, not a survey of deployed products.

\section{Background and related work}
CXF can encode FIDO2 HMAC state~\cite{cxf}. The CXP working draft defines a
transport profile~\cite{cxp}; our single-document CXP/HPKE binding is an
experimental project implementation and does not cover its full ZIP/JWE
wire format. HPKE follows RFC~9180~\cite{hpke}. A WebAuthn PRF result is
returned through client extension results, outside the ordinary signed
assertion bytes~\cite{webauthn}. Consequently, signature and PRF oracles
answer separate questions. Ramat et al. compare passkey user experiences
across 111 websites~\cite{ramat}; Jannett et al. measure deployment and
security on passkey-enabled websites~\cite{jannett}. B\"uttner and
Gruschka compare the security and availability of device-bound and synced
credentials~\cite{buttner}. Daffalla et al. study passkey abuse under
interpersonal threat models~\cite{daffalla}, while Islam et al. propose
detection of stolen cloud-stored passkeys~\cite{islam}. Henarejos proposes
authenticated export and restoration of device-bound credentials and a
falsifiable evaluation plan~\cite{henarejos}. These studies respectively
observe relying parties, user-facing management, cloud compromise, or a
proposed recovery design. None of those reported measurements establishes
whether a controlled credential retains its extension behavior after a
specific provider-to-provider import. That is the narrower question here.

Our unit is an attempted migration of a controlled credential, with
separate post-import oracles for signing, extension behavior, and retry
state. The method makes loss of observability an explicit outcome, rather
than interpreting a successful assertion as evidence about unexecuted
properties. The controlled routes expose that distinction, while the
selected provider cases test only a few concrete code paths. We do not
infer website adoption, attack prevalence, or provider prevalence.

\section{Study design}
\subsection{Units, outcomes, and exclusions}
The v1.5 protocol crossed 256 synthetic ES256 credentials (32 in each of
eight strata) with 12 routes and two
repetitions: 6,144 attempts in 96 route-by-stratum groups. The strata
covered basic credentials, PRF with and without UV, largeBlob, credBlob,
a payment-format marker, combined extensions, and a future optional
member. They were constructed controls, not random draws. A route
could reject the document before WebAuthn. For each applicable property we
record PASS, FAIL, NOT\_APPLICABLE, or NOT\_EVALUABLE. The last outcome is
needed where CDP cannot inject PRF or credBlob state.

One oracle checks CXF representation and policy. Each imported attempt
performs two WebAuthn assertions: the first binds a challenge to the attempt,
credential, route, repetition, public key, user handle, RP ID, and origin;
the second binds its challenge to the first and to a canonical digest of
extension observations. A separate verifier checks both ES256 signatures,
the RP hash, user-presence and verification flags, and counter. Extension
oracles examine executed behavior where the
interface permits it; representation loss alone is never relabeled as an
executed PRF failure. A deterministic in-memory simulation tests transfer
atomicity and retry behavior. These oracles are not interchangeable.
Campaign percentages are design-weighted proportions, without population
confidence intervals; paired repetitions are stability checks, not
independent providers. The browser and provider
follow-ups were designed after inspecting earlier results and are
exploratory rather than preregistered confirmations.

\subsection{Threat and implementation boundary}
Two separately written project implementations exchange one CXF document
in Python-to-Node and Node-to-Python directions and reject a modified
ciphertext. This tests that experimental binding. Exporter identity,
authorization, durable provider storage, and the public CXP ZIP/JWE
envelope are outside the tested profile. The transaction model is in
memory and makes no claim about crash durability.

The reference route models an importer that retains supported fields.
The strict route refuses a credential when it cannot preserve a required
property. The compatible-lossy route accepts a credential while dropping
selected extension state; the legacy route has a different loss policy.
These are controlled policies written for this study, not named vendors.
Direct, round-trip, and multihop chains expose whether a later capable
destination can recover a field already lost at an earlier hop.

The source credential, expected CXF fields, and route policy are known
to the harness. Its representation oracle therefore has ground truth that
a real RP usually lacks. The browser oracle sees only what a WebAuthn
ceremony and client extension results expose. This separation prevents a
missing serialized seed from being counted as a failed browser PRF
evaluation. The same distinction applies to a payment-format marker
without a Secure Payment Confirmation ceremony.

\section{Synthetic campaign results}
Of 6,144 attempts, 1,024 were rejected before ceremony and 5,120 were
imported. Every executed assertion among the imported attempts was
accepted. Table~\ref{tab:c1} separates those accepted assertions from other
oracles. The complete observable PASS yield was 2,304/6,144 (37.50\%).
This describes the balanced test design, not the probability that an
arbitrary user or product will succeed.

\begin{table}[t]
\centering
\caption{C1 synthetic attempts; denominator 6,144.}
\label{tab:c1}
\begin{tabular}{lrr}
\toprule
Classification & Count & Share \\
\midrule
Complete observable PASS & 2,304 & 37.50\% \\
Visible degradation & 960 & 15.62\% \\
Silent degradation & 1,344 & 21.88\% \\
Not evaluable & 512 & 8.33\% \\
Pre-ceremony rejection & 1,024 & 16.67\% \\
\bottomrule
\end{tabular}
\end{table}

\begin{table*}[t]
\centering
\caption{C1 outcomes by route. Each route has 512 designed attempts.
P: complete observable PASS; V: visible degradation; S: silent
degradation; NE: not evaluable. Counts are exact for this design.}
\label{tab:routes}
\begin{tabular}{llrrrrrr}
\toprule
Route & Chain & Imported & Rejected & P & V & S & NE \\
\midrule
R01 & ref--strict & 512 & 0 & 256 & 0 & 0 & 256 \\
R02 & ref--lossy & 512 & 0 & 256 & 256 & 0 & 0 \\
R03 & ref--legacy & 512 & 0 & 128 & 0 & 384 & 0 \\
R04 & ref--strict--ref & 512 & 0 & 256 & 0 & 0 & 256 \\
R05 & ref--lossy--ref & 512 & 0 & 256 & 256 & 0 & 0 \\
R06 & ref--legacy--ref & 512 & 0 & 128 & 0 & 384 & 0 \\
R07 & ref--strict--lossy & 512 & 0 & 256 & 256 & 0 & 0 \\
R08 & ref--lossy--strict & 256 & 256 & 256 & 0 & 0 & 0 \\
R09 & ref--strict--legacy & 512 & 0 & 128 & 0 & 384 & 0 \\
R10 & ref--legacy--strict & 128 & 384 & 128 & 0 & 0 & 0 \\
R11 & ref--lossy--legacy--strict & 128 & 384 & 128 & 0 & 0 & 0 \\
R12 & ref--strict--lossy--legacy--ref & 512 & 0 & 128 & 192 & 192 & 0 \\
\midrule
Total & & 5,120 & 1,024 & 2,304 & 960 & 1,344 & 512 \\
\bottomrule
\end{tabular}
\end{table*}

\begin{table*}[t]
\centering
\caption{C1 outcomes by constructed credential stratum.
Each stratum has 768 designed attempts over all routes and repetitions.
Labels F0--F7 are defined in the text.}
\label{tab:strata}
\begin{tabular}{lrrrrrr}
\toprule
Stratum & Imported & Rejected & P & V & S & NE \\
\midrule
F0 basic & 768 & 0 & 768 & 0 & 0 & 0 \\
F1 PRF with UV & 576 & 192 & 0 & 256 & 192 & 128 \\
F2 PRF without UV & 576 & 192 & 0 & 256 & 192 & 128 \\
F3 largeBlob & 640 & 128 & 384 & 0 & 256 & 0 \\
F4 credBlob & 576 & 192 & 0 & 256 & 192 & 128 \\
F5 payment marker & 640 & 128 & 384 & 0 & 256 & 0 \\
F6 combined extensions & 576 & 192 & 0 & 192 & 256 & 128 \\
F7 optional member & 768 & 0 & 768 & 0 & 0 & 0 \\
\midrule
Total & 5,120 & 1,024 & 2,304 & 960 & 1,344 & 512 \\
\bottomrule
\end{tabular}
\end{table*}

Among imported attempts, 512 (10.00\%) combined valid login with failure
of an executed UV or largeBlob browser oracle. PRF or credBlob
representation loss occurred in 1,792 (35.00\%); this category overlaps
the former and does not establish an executed PRF failure in those cells.
The union of non-payment failures was 2,048/5,120 (40.00\%). A
payment-format marker raises the design sensitivity to 2,304/5,120
(45.00\%), but is not a Secure Payment Confirmation ceremony. C2 malformed
document controls comprised 80 cases across ten families and eight strata;
64 were rejected by structural validation, duplicate policy, or strict
preservation. C3 exercised 960 in-memory failure sequences and 1,920 events:
832 rolled back and converged to one copy on retry, while the 128 deliberate
legacy late-failure controls retained a provisional and duplicated on retry.
Neither C2 nor C3 is included in the C1 percentages.

Table~\ref{tab:routes} shows the route dependence directly. R03, R06,
and R09 each contain 384 silent-degradation attempts, despite 512
imports and valid assertions. In the paired R02/R05 comparison, returning
to the reference provider leaves the same 256 visible-degradation
attempts: a later capable endpoint cannot reconstruct dropped source
state. Strict handling after a lossy hop rejects 256 attempts in R08;
after a legacy hop it rejects 384 in R10. These are consequences of the
specified control policies, not failure rates for real providers.

Table~\ref{tab:strata} separates which controls generate the aggregate.
The basic and optional-member strata each have 768 observable PASS
attempts. PRF-bearing and credBlob strata have no complete PASS because
their positive browser behavior cannot be established through the CDP
injection path even when their serialized representation is retained.
The 512 false-reassurance cases occur in largeBlob-bearing strata F3
and F6: 256 in each, with accepted assertions and a failing executed
browser behavior oracle. These 512 cases are a subset of the broader
semantic-degradation counts, not an additional category to add to
Table~\ref{tab:c1}.

The C2 family outcome matters more than its pooled rejection count.
Eight unknown-optional-member and eight minor-version-increase cases
were imported but assigned NOT\_EVALUABLE; a credential-ID collision
was rejected in eight cases while its normative class remained
ambiguous. C3's 128 failed atomicity and idempotence controls were
exactly the two late failure points of the legacy model, 64 each.
Because C3 has no process crash or durable storage, its result cannot
predict recovery behavior of a password manager.

\section{Executable follow-up controls}
\subsection{Experimental exchange and browser PRF loss}
Python and Node recovered identical CXF documents in both directions and
rejected ciphertext alteration. Software PRF evaluators agreed on retained
CXF HMAC seeds and detected their removal. This is a functional software
control, not a browser ceremony.

A Chromium virtual authenticator created a PRF-enabled passkey. Three
source assertions showed a PRF result stable for one salt and different
for a second. The runner exported the signing credential via CDP, carried
the CXF document through the experimental exchange, and injected the
credential into a new virtual authenticator. A fourth assertion
authenticated at the destination, while the PRF result for the original
salt was absent. CDP did not expose a PRF seed. This is a loss control for
that test path, not a finding about a commercial credential provider.

A Windows 2025 GitHub Actions replay recorded six exchange and nine
browser checks as true. A second replay retained four public signatures,
the registration public key, challenges, authenticator and client data,
and SHA-256 digests of observed PRF outputs. A separate verifier
checked four ES256 signatures, RP binding, challenge uniqueness, UV
flags, credential identity, and four PRF digest relationships. The archive
contains no private key, seed, or raw PRF output. Comparing hashes proves
relations among observed outputs, not the HMAC computation from a seed
that CDP did not reveal.

\subsection{Public-draft envelope structure probe}
We added a synthetic shape probe for the CXP Working Draft's visible
request and response members, offered parameter selection, and the
illustrated ZIP layout with an index JWE and a document JWE.
The positive fixture passes those structural checks. Three negative
responses with a version mismatch, unoffered archive, or malformed
payload encoding are rejected. Additional tests reject unexpected ZIP
paths and a plain JSON file substituted for a JWE. The fixture uses
standard direct-key AES-GCM JWE solely to exercise the container shape.
It does not claim that this maps HPKE to JWE. The draft does not give a
complete HPKE-to-JWE binding or include its narrative signed-challenge
fields in the displayed request and response objects. Consequently
JWE authentication under negotiated HPKE, challenge verification,
CXF plaintext recovery, and provider interoperability remain
NOT\_EVALUABLE in this probe~\cite{cxp}. It is a reproducible parsing
control, not a CXP conformance or security result.

To identify an external implementation target, we compared the fixture's
request and response fields with a pinned Bitwarden CXP Rust type
definition~\cite{bitwardencxp}. That source declares neither the draft's
\texttt{archive} and \texttt{mode} request fields nor its \texttt{archive}
response field, and stores the version as an eight-bit value rather than
the draft's 16-bit type. The comparison is executable from a frozen
field snapshot. We compiled the pinned crate with locked dependencies in
CI and parsed synthetic request and response objects. Both parsed, but
serialization dropped the request's \texttt{archive} and \texttt{mode}
fields and the response's \texttt{archive} field. A draft-valid version
of 256 was rejected. This is a type-parser round-trip gap, not a provider
transfer or encrypted-payload result; it cannot establish interoperability.

\subsection{Selected provider case: Fenko}
We selected Fenko Vault 0.12.0~\cite{fenko} after inspecting its CXF
path. Two fresh profiles per repetition imported synthetic basic and
PRF-bearing credentials. The credentials signed valid assertions. A
public verifier rechecked 16 signatures and six PRF observation bindings.
The source CXF HMAC seeds were not retained in the observed import/export
path; the browser PRF result differed from an oracle derived from those
source seeds.

A same-provider round trip between two Fenko profiles produced matching
PRF outputs in two selected repetitions. A separate verifier recalculated
24 signatures and 12 PRF bindings. This continuity does not imply that
the original seeds survived: this code path derived a value from the
imported signing key. Functional equality and preservation of the
represented secret are separate claims. Two repetitions are not a
product-level reliability estimate.

\subsection{Positive software control: Proton}
The Proton Pass common 2.0.0 software authenticator~\cite{proton}
provided a selected positive control. In two exploratory repetitions,
the runner imported CXF, exported it, and reimported into a second
software authenticator. Both HMAC seeds survived. With user verification,
the expected PRF was reproduced for two salts before and after reimport.
A modified-seed control changed the result while the credential still
signed. The public verifier rechecked 12 signatures and 12 PRF digest
comparisons, including four modified-seed controls. The isolated report
cannot reconstruct the HMAC from a withheld seed; that requires repeating
the pinned software run. This route used pinned WASM, not a
browser passkey provider or the installed Proton Pass application; no
no-UV ceremony was run.

\section{Discussion and limitations}
The browser loss control and Fenko case show different ways a signing
credential can outlive its original PRF state. The Proton software route
shows that preserving CXF seeds and reproducing the expected PRF is
possible in one selected implementation. These cases are not a
representative sample of provider behavior.

The public CXP wire format was not exchanged between two independent
providers. A human-reported Android Bitwarden-to-Google transfer lacks an
instrumented destination assertion and transfer receipt and is excluded
from measured results. CDP cannot inject original PRF state; positive
browser-to-browser restoration of CXF HMAC seeds is NOT\_EVALUABLE. The
Proton module result cannot fill that gap. The transaction simulation
cannot support claims about storage durability, and the synthetic campaign
cannot support product prevalence or superiority claims.

\section{Conclusion}
PasskeyTransit uses distinct oracles for identity, signing, extensions,
and transaction semantics, marking properties an interface cannot expose.
The synthetic campaign and selected controls show valid authentication
alongside changed or missing properties. Public replays verify observed
signatures; provider interoperability and mobile transfer remain open.

\appendix
\section{Open Science}
The project archive contains source, locked dependencies, protocols,
synthetic reports, replay ZIPs, and verifiers. A clean extraction passed
on the same team's Windows host on 20 October 2026; this is not
independent replication. A staged anonymous archive includes synthetic
C1/C2/C3 rows, analysis tables, the structural CXP probe, the executed
Rust type-parser report, the original Phase 14 browser replay ZIP, and a
local Fenko/Proton replay bundle. Verifiers recompute 96 C1 cells, check
member hashes and public signatures, and match the Phase 14 reports byte
for byte. GitHub run metadata stays in a private sidecar. The provider
replay has exact runner snapshots but no external CI attestation. Neither
archive contains vaults, private keys, original PRF seeds, or raw PRF
outputs. A stable anonymous review URL remains to be prepared.

\section{Ethical considerations}
All reported software credentials and RP accounts were synthetic; the
browser control used virtual authenticators and a localhost RP. No real
passkey or private vault is distributed. The excluded Android observation
is not independently verified. Provider findings apply to tested routes
and revisions only.

\begin{thebibliography}{16}
\bibitem{cxf} FIDO Alliance. Credential Exchange Format v1.0 Proposed Standard Errata, 2026. \url{https://fidoalliance.org/specs/cx/cxf-v1.0-ps-errata-20260309.html}.
\bibitem{cxp} FIDO Alliance. Credential Exchange Protocol v1.0 Working Draft, 3 October 2024. \url{https://fidoalliance.org/specs/cx/cxp-v1.0-wd-20241003.html}.
\bibitem{hpke} R. Barnes et al. Hybrid Public Key Encryption. RFC 9180, 2022. \url{https://www.rfc-editor.org/rfc/rfc9180}.
\bibitem{webauthn} W3C. Web Authentication Level 3, Recommendation, 25 August 2026. \url{https://www.w3.org/TR/2026/REC-webauthn-3-20260825/}.
\bibitem{ramat} B. Ramat, D. Kartchner, M. Clark, and K. Seamons. Passkeys in the Wild: A Systematic Study of FIDO2 User Experience Consistency Across Websites. SOUPS, pp. 21--39, 2026. \url{https://www.usenix.org/conference/soups2026/presentation/ramat}.
\bibitem{jannett} L. Jannett, A. Mayer, M. Westers, V. Mladenov, C. Mainka, and J. Schwenk. The State of Passkeys: Studying the Adoption and Security of Passkeys on the Web. USENIX Security, pp. 5039--5058, 2026. \url{https://www.usenix.org/conference/usenixsecurity26/presentation/jannett}.
\bibitem{buttner} A. B{\"u}ttner and N. Gruschka. Device-Bound vs. Synced Credentials: A Comparative Evaluation of Passkey Authentication. ICISSP, 2025. \url{https://arxiv.org/abs/2501.07380}.
\bibitem{daffalla} A. Daffalla et al. A Framework for Abusability Analysis: The Case of Passkeys in Interpersonal Threat Models. USENIX Security, pp. 7819--7838, 2025. \url{https://www.usenix.org/conference/usenixsecurity25/presentation/daffalla}.
\bibitem{islam} M. Islam et al. Detecting Compromise of Passkey Storage on the Cloud. USENIX Security, pp. 7743--7762, 2025. \url{https://www.usenix.org/conference/usenixsecurity25/presentation/islam}.
\bibitem{henarejos} P. Henarejos. Vaulted Passkeys: A Device-Bound Proposal for Authenticated Credential Export and Import. arXiv:2608.13806, 2026. \url{https://arxiv.org/abs/2608.13806}.
\bibitem{bitwardencxp} Bitwarden. Credential Exchange Protocol Rust types, pinned commit \texttt{32da4c8b0a4abdd0b46fae79c31e0fdb50150396}, 2026. \url{https://github.com/bitwarden/credential-exchange/blob/32da4c8b0a4abdd0b46fae79c31e0fdb50150396/credential-exchange-protocol/src/lib.rs}.
\bibitem{fenko} FenkoHQ. Passkey Vault 0.12.0, pinned commit \texttt{2d873fbe25257250de6e322e9d785b97b17a544f}. \url{https://github.com/FenkoHQ/passkey-vault}.
\bibitem{proton} Proton. Proton Pass common 2.0.0, pinned commit \texttt{69a0927f0fe7071243d6c38e404459a69210fdd4}. \url{https://github.com/protonpass/proton-pass-common}.
\end{thebibliography}
\end{document}
